\documentclass[journal]{IEEEtran}

\usepackage[T1]{fontenc}
\usepackage{amsmath,amssymb,amsfonts,amsthm}
\usepackage{booktabs}
\usepackage{tikz}
\usetikzlibrary{positioning,arrows.meta,fit,calc}
\usepackage{multirow}
\usepackage{xcolor}
\usepackage{bm}
\usepackage{algorithm}
\usepackage{algorithmic}
\usepackage{graphicx}
\usepackage{url}
\usepackage{cite}
\usepackage{balance}
\usepackage{subcaption}

\newtheorem{proposition}{Proposition}
\newtheorem{corollary}{Corollary}
\newtheorem{definition}{Definition}
\newtheorem{remark}{Remark}

\definecolor{dyncolor}{RGB}{66,133,244}
\definecolor{fixcolor}{RGB}{52,168,83}
\definecolor{newcolor}{RGB}{180,120,0}
\definecolor{tempcolor}{RGB}{171,71,188}
\definecolor{tcncolor}{RGB}{220,38,38}

\title{NC-Conv: Noise-Conditioned Dual-Path Convolution and Temporal State Space Models for Degradation-Robust Vision}

\author{Jin~Ho~Choi,~\IEEEmembership{Member,~IEEE}%
\thanks{J.~H.~Choi is the CEO of SmartEar Co., Ltd., 1143-1 Geumri, Yuga-eup, Dalseong-gun, Daegu 43020, South Korea (e-mail: jinhochoi@smartear.co.kr).}%
\thanks{This paper extends the authors' conference work presented at ACCV 2026. Extensions include: (1)~NC-TCN Vision with dilated multi-scale convolution (Section~\ref{sec:nctcn}), (2)~formal noise propagation proofs (Section~\ref{sec:theory}), (3)~comprehensive CIFAR-10-C evaluation with 19 corruptions (Section~\ref{sec:cifar10c}), (4)~CULane real-world lane detection (Section~\ref{sec:culane}), (5)~NC-VLU ultra-lightweight vision-language model (Section~\ref{sec:ncvlu}), and (6)~MCU deployment analysis (Section~\ref{sec:deploy}).}%
\thanks{Manuscript submitted 2027.}}

\begin{document}
\maketitle

% ============================================================
% ABSTRACT
% ============================================================
\begin{abstract}
We present a family of \textbf{noise-conditioned (NC)} vision architectures that transfer audio noise-conditioning theory to visual degradation robustness.
The core insight is that input-dependent operations (dynamic convolution, channel gating) amplify noise as $O(\sigma_n^{4+})$, while fixed-weight operations maintain $O(\sigma_n^2)$---mirroring the $O(\sigma_n^6)$ vs.\ $O(\sigma_n^2)$ gap in selective vs.\ fixed state space models established in our audio NC-SSM work.
Each NC block blends a \textcolor{fixcolor}{\emph{static}} path (robust, $O(\sigma_n^2)$) with a \textcolor{dyncolor}{\emph{dynamic}} path (expressive, $O(\sigma_n^{4+})$) via a learned quality gate $\sigma$ that automatically adapts to input degradation level.

We propose three complementary components:
(1)~\textbf{NC-Conv}, a dual-path convolutional block achieving $+5.7\%$ corrupted accuracy on CIFAR-10-C;
(2)~\textbf{NC-TCN Vision}, extending our audio NC-TCN with dilated multi-scale convolution (dilation $\{1,2,4\}$) and per-spatial $\sigma$, achieving $\mathbf{+8.1\%}$ with fog $+18.2\%$;
(3)~\textbf{Bidirectional Temporal NC-SSM} with per-frame quality gating for video, recovering $\mathbf{+35.2\%}$ on the darkest tunnel frame while preserving clean accuracy within $1\%$.

On CULane real-world lane detection, NC-Conv achieves $+21.7\%$ on fog conditions.
The quality gate learns degradation profiles \emph{without explicit quality labels}, purely through task loss.
At 253K parameters (253\,KB INT8), the system deploys on a \$8 Cortex-M7 MCU at 28\,FPS, establishing noise-conditioned computation as a modality-agnostic framework for degradation-robust edge perception.
\end{abstract}

\begin{IEEEkeywords}
Corruption robustness, noise conditioning, dynamic convolution, lane detection, state space models, vision-language models, edge AI, MCU deployment
\end{IEEEkeywords}


% ============================================================
% I. INTRODUCTION
% ============================================================
\section{Introduction}
\label{sec:intro}

\IEEEPARstart{D}{eep} vision models deployed in safety-critical applications---autonomous driving, surveillance, and industrial inspection---must maintain accuracy under adverse conditions including fog, low light, sensor noise, and compression artifacts.
These degradation scenarios arise precisely when robust perception is most critical: a tunnel entrance creates extreme dynamic range transitions, night driving reduces visibility by orders of magnitude, and weather conditions introduce spatially varying corruption.

While corruption augmentation during training~\cite{hendrycks2020augmix,hendrycks2021deepaugment} improves average robustness, it treats all inputs identically regardless of their quality, missing the opportunity to adapt processing to the specific degradation level of each input.
Dynamic convolution~\cite{chen2020dynamic} and channel attention~\cite{hu2018squeeze} increase model expressiveness through input-dependent parameterization, but our analysis reveals that this very input-dependence creates \emph{multiplicative noise amplification}---the same vulnerability identified in selective state space models (SSMs) for audio~\cite{choi2026ncssm}.

\subsection{From Audio Noise to Visual Degradation}

In our prior work on noise-robust keyword spotting~\cite{choi2026ncssm}, we established that selective SSMs exhibit inherent noise robustness through input-dependent dynamics, but suffer from multiplicative noise propagation under severe conditions:
\begin{equation}
\text{Var}[y_{\text{selective}}] = O(\sigma_n^6), \quad
\text{Var}[y_{\text{fixed}}] = O(\sigma_n^2).
\label{eq:audio_var}
\end{equation}
The NC-SSM solution---blending selective and fixed dynamics based on local signal quality (SNR)---transfers to vision when we observe that:
\begin{itemize}
\setlength\itemsep{0pt}
\item \textbf{Images are not sequential data}: unlike audio, spatial patches lack natural ordering, making CNNs the appropriate spatial backbone rather than SSMs.
\item \textbf{Video is sequential}: the temporal frame axis \emph{is} naturally sequential, making SSMs appropriate for temporal modeling.
\item \textbf{Input-dependent CNN operations amplify noise}: channel gating ($\alpha(x) \odot \text{Conv}(x)$) creates $O(\sigma_n^{4+})$ variance---structurally analogous to the SSM's $O(\sigma_n^6)$.
\end{itemize}

\subsection{Contributions}

This paper presents a comprehensive noise-conditioned vision framework:

\begin{enumerate}
\setlength\itemsep{0pt}
\item \textbf{Noise propagation theory for CNNs}: We prove that input-dependent channel modulation creates $O(\sigma_n^{4+})$ output variance, and that $\sigma$-gated interpolation with a fixed path reduces this to $O(\sigma_n^{2+\epsilon})$ (Section~\ref{sec:theory}).

\item \textbf{NC-Conv}: A dual-path convolutional block with learned quality gate $\sigma$, achieving $+5.7\%$ on corrupted CIFAR-10 and $+21.7\%$ on CULane fog (Section~\ref{sec:ncconv}).

\item \textbf{NC-TCN Vision}: Extending our audio NC-TCN~\cite{choi2026nctcn} with dilated multi-scale convolution (dilation $\{1,2,4\}$, receptive field $15\times 15$) and per-spatial $\sigma$, achieving $+8.1\%$ average corrupted accuracy (Section~\ref{sec:nctcn}).

\item \textbf{Bidirectional Temporal NC-SSM}: Per-frame quality-gated temporal context transfer for video, recovering $+35.2\%$ on the darkest tunnel frame with clean accuracy preserved within $1\%$ (Section~\ref{sec:temporal}).

\item \textbf{Severity scaling}: NC-Conv's advantage grows monotonically with corruption severity ($+0.5\%$ to $+2.9\%$), directly validating the noise propagation theory (Section~\ref{sec:cifar10c}).

\item \textbf{Implicit quality learning}: The $\sigma$ gate and temporal quality gate learn degradation profiles without explicit quality labels, purely through task loss (Section~\ref{sec:analysis}).

\item \textbf{MCU deployment}: At 253K parameters (253\,KB INT8), the system deploys on a \$8 Cortex-M7 MCU at 28\,FPS with 35\,ms latency and 0.5\,W power (Section~\ref{sec:deploy}).
\end{enumerate}


% ============================================================
% II. RELATED WORK
% ============================================================
\section{Related Work}
\label{sec:related}

\subsection{Corruption Robustness}
The CIFAR-10-C and ImageNet-C benchmarks~\cite{hendrycks2019benchmarking} standardized corruption robustness evaluation with 19 corruption types across 5 severity levels.
AugMix~\cite{hendrycks2020augmix} improves robustness through diverse augmentation mixing.
DeepAugment~\cite{hendrycks2021deepaugment} uses neural-network-based augmentation.
PixMix~\cite{pixmix2022} combines natural and synthetic corruptions.
All these methods apply uniform processing regardless of input quality.
NC-Conv uniquely combines augmentation with \emph{quality-adaptive} dual-path routing, where the processing itself changes based on estimated input degradation.

\subsection{Dynamic and Conditional Computation}
Dynamic convolution~\cite{chen2020dynamic} uses input-dependent kernel attention over multiple basis kernels.
SE-Net~\cite{hu2018squeeze} provides squeeze-and-excitation channel recalibration.
CondConv~\cite{condconv2019} parameterizes convolution conditioned on input.
Mixture-of-Experts (MoE) architectures~\cite{shazeer2017outrageously} route inputs through specialized sub-networks.
NC-Conv differs by grounding dual-path routing in \emph{noise propagation theory}, providing a principled static fallback under degradation rather than maximizing average expressiveness.

\subsection{Low-Light and Adverse Vision}
Retinex-based methods~\cite{land1977retinex,wei2018deepretinex} decompose illumination and reflectance.
Zero-DCE~\cite{guo2020zerodce} performs zero-reference curve estimation for low-light enhancement.
Image-Adaptive YOLO (IA-YOLO)~\cite{iayolo2022} integrates image enhancement with object detection.
PE-YOLO~\cite{peyolo2023} proposes a parameter-efficient enhancement-detection pipeline.
These treat enhancement as preprocessing, adding parameters and latency.
Our NC framework integrates quality-adaptive processing \emph{within} the backbone at near-zero parameter cost.

\subsection{Lane Detection}
Ultra-Fast Lane Detection (UFLD)~\cite{qin2020ultrafast} formulates lanes as row-anchor classification.
CLRNet~\cite{zheng2022clrnet} uses cross-layer refinement.
CLRerNet~\cite{clrernet2024} improves confidence with LaneIoU.
These methods require 0.9--9M parameters and CPU/GPU hardware.
Our NC-Conv lane detector operates at 253K parameters on a \$8 MCU.

\subsection{Noise-Conditioned Models}
Our audio NC-SSM~\cite{choi2026ncssm} blends selective ($O(\sigma_n^6)$) and LTI ($O(\sigma_n^2)$) SSM dynamics per sub-band based on SNR.
Our audio NC-TCN~\cite{choi2026nctcn} replaces sequential SSM with dilated causal convolution for parallelism and INT8 stability.
This work extends both frameworks to vision (spatial) and video (temporal).

\subsection{Vision-Language Models}
CLIP~\cite{radford2021clip} learns visual representations from natural language supervision at scale (400M image-text pairs, 428M parameters).
MobileCLIP~\cite{mobileclip2024} reduces to 20M parameters.
TinyCLIP~\cite{tinyclip2024} further compresses to 40M.
Our NC-VLU targets MCU-scale visual keyword spotting at $<$1M parameters---100$\times$ smaller than the smallest existing VLM.


% ============================================================
% III. NOISE PROPAGATION THEORY
% ============================================================
\section{Noise Propagation in CNNs}
\label{sec:theory}

We establish the noise propagation properties of standard CNN operations, extending the audio SSM analysis~\cite{choi2026ncssm} to spatial convolutions.

\begin{definition}[Static and Dynamic Convolution]
\label{def:conv}
A \textbf{static} convolution computes $h_s = W * x$ with fixed kernel $W$.
A \textbf{dynamic} convolution computes $h_d = (W * x) \odot \alpha(x)$, where $\alpha(x) = \text{Sigmoid}(\text{FC}(\text{GAP}(x))) \in \mathbb{R}^C$ is a channel-wise input-dependent gate.
\end{definition}

\begin{proposition}[Static Noise Variance]
\label{prop:static}
For noisy input $x = s + n$, $n \sim \mathcal{N}(0, \sigma_n^2 I)$, and fixed kernel $W$:
\begin{equation}
\text{Var}[h_s] = \|W\|_F^2 \cdot \sigma_n^2 = O(\sigma_n^2).
\end{equation}
\end{proposition}

\begin{proof}
Since $W$ is input-independent, $h_s = W * s + W * n$, and $\text{Var}[W * n] = \|W\|_F^2 \sigma_n^2$.
\end{proof}

\begin{proposition}[Dynamic Noise Amplification]
\label{prop:dynamic}
For the dynamic path $h_d = (W * x) \odot \alpha(x)$ with $x = s + n$:
\begin{align}
\text{Var}[h_d] &= \underbrace{\bar{\alpha}^2 \|W\|_F^2 \sigma_n^2}_{O(\sigma_n^2)} + \underbrace{(W*s)^2 \sigma_\alpha^2}_{O(\sigma_n^2 \sigma_\alpha^2)} \nonumber \\
&\quad + \underbrace{\|W\|_F^2 \sigma_n^2 \sigma_\alpha^2}_{O(\sigma_n^2 \sigma_\alpha^2)} + \underbrace{O(\sigma_n^4)}_{\text{high-order}},
\label{eq:dynamic_var}
\end{align}
where $\bar{\alpha} = \mathbb{E}[\alpha(x)]$ and $\sigma_\alpha^2 = \text{Var}[\alpha(x)]$.
Since $\alpha$ depends on $x$ through a nonlinear projection:
\begin{equation}
\sigma_\alpha^2 = \text{Var}[\alpha(s+n)] \approx \|\nabla_x \alpha\|^2 \sigma_n^2 + O(\sigma_n^4),
\end{equation}
yielding $\text{Var}[h_d] = O(\sigma_n^2) + O(\sigma_n^4) + O(\sigma_n^6)$.
\end{proposition}

\begin{proposition}[NC-Conv Variance Bound]
\label{prop:nc}
The NC-Conv interpolation $h = \sigma \cdot h_d + (1-\sigma) \cdot h_s$ satisfies:
\begin{equation}
\text{Var}[h] = \sigma^2 \text{Var}[h_d] + (1{-}\sigma)^2 \text{Var}[h_s] + 2\sigma(1{-}\sigma)\text{Cov}[h_d, h_s].
\end{equation}
At $\sigma=0$: $\text{Var}[h] = O(\sigma_n^2)$ (maximum robustness).
At $\sigma=1$: $\text{Var}[h] = O(\sigma_n^{4+})$ (maximum expressiveness).
\end{proposition}

\begin{corollary}[Optimal $\sigma$]
\label{cor:optimal}
The optimal gate $\sigma^*(\sigma_n) \to 0$ as noise increases and $\sigma^* \to 1$ as noise vanishes, producing the monotonic severity scaling observed experimentally (Table~\ref{tab:severity}).
\end{corollary}

\begin{proposition}[Superset Guarantee]
\label{prop:superset}
$\mathcal{H}_{\text{NC}} \supseteq \mathcal{H}_{\text{std}}$: setting $\sigma=1\;\forall x$ recovers the standard dynamic network. NC-Conv can only improve, never degrade, performance.
\end{proposition}

\begin{table}[t]
\centering
\caption{Noise variance comparison across modalities. Input-dependent operations amplify noise super-linearly in both audio SSMs and vision CNNs.}
\label{tab:noise_theory}
\small
\begin{tabular}{@{}lcc@{}}
\toprule
\textbf{Operation} & \textbf{Audio (SSM)} & \textbf{Vision (CNN)} \\
\midrule
Fixed & LTI: $O(\sigma_n^2)$ & Static Conv: $O(\sigma_n^2)$ \\
Input-dependent & Selective: $O(\sigma_n^6)$ & Dynamic: $O(\sigma_n^{4+})$ \\
\textbf{NC (ours)} & $O(\sigma_n^{2+\epsilon})$ & $O(\sigma_n^{2+\epsilon})$ \\
\bottomrule
\end{tabular}
\end{table}


% ============================================================
% IV. NC-CONV ARCHITECTURE
% ============================================================
\section{NC-Conv: Noise-Conditioned Convolution}
\label{sec:ncconv}

\subsection{NC-Conv Block}

Each NC-Conv block contains two convolutional paths and a quality gate:

\paragraph{Static path (robust, $O(\sigma_n^2)$)}
A depthwise separable convolution with fixed weights:
\begin{equation}
h_s = \text{BN}(\text{DWConv}_{\text{static}}(x)).
\end{equation}

\paragraph{Dynamic path (expressive, $O(\sigma_n^{4+})$)}
A depthwise convolution with SE-style~\cite{hu2018squeeze} channel modulation:
\begin{equation}
h_d = \text{BN}(\text{DWConv}_{\text{dyn}}(x)) \odot \alpha(x),
\end{equation}
where $\alpha(x) = \text{Sigmoid}(\text{FC}(\text{GAP}(x)))$.

\paragraph{Quality gate $\sigma$}
A lightweight MLP estimates per-sample quality:
\begin{equation}
\sigma = \text{Sigmoid}(\text{FC}(\text{SiLU}(\text{FC}(\text{GAP}(x))))).
\end{equation}

\paragraph{NC blend}
\begin{equation}
h = \sigma \cdot h_d + (1-\sigma) \cdot h_s,
\label{eq:nc_blend}
\end{equation}
followed by pointwise convolution and residual connection.

\subsection{Per-Spatial $\sigma$ Extension}
\label{sec:perspatial}

The per-sample $\sigma$ gate (Eq.~\ref{eq:nc_blend}) uses global average pooling, which cannot detect spatially localized corruption (e.g., impulse noise affecting individual pixels, partial shadow).
The per-spatial extension computes $\sigma \in \mathbb{R}^{B \times 1 \times H \times W}$ using a small convolutional network:
\begin{equation}
\sigma(x) = \text{Sigmoid}(\text{Conv}_{1\times 1}(\text{SiLU}(\text{DWConv}_{3\times 3}(\text{Conv}_{1\times 1}(x))))),
\end{equation}
enabling different blend ratios at each spatial location---the direct analog of audio NC-SSM's per-sub-band $\sigma$.

\subsection{Training}

NC-Conv uses corruption augmentation (30\% of mini-batch randomly corrupted with one of five types at random severity 1--3).
The $\sigma$ gate learns quality routing \emph{implicitly} through the classification loss.
Explicit $\sigma$ supervision with binary labels (clean$\to$1, noisy$\to$0) degrades performance by $-4.4\%$, as the classification loss provides a better, softer training signal.


% ============================================================
% V. NC-TCN VISION
% ============================================================
\section{NC-TCN Vision: Dilated Multi-Scale Extension}
\label{sec:nctcn}

Our audio NC-TCN~\cite{choi2026nctcn} replaces sequential SSM with dilated causal convolution for full parallelism and INT8 quantization stability.
We apply the same principle to vision: each NC-TCN Vision block uses \textbf{dilated depthwise Conv2d} with dilation schedule $\{1, 2, 4\}$, combined with per-spatial $\sigma$ and gated activation:
\begin{equation}
h_d = \text{SiLU}(\text{DWConv}_{d}^{\text{dyn}}(x) + g_x) \odot \sigma_z(g_z),
\end{equation}
where $g_x, g_z = \text{split}(\text{Conv}_{1\times 1}(x))$ and dilation $d$ expands the receptive field.
The effective receptive field across three stacked blocks is:
\begin{equation}
\text{RF} = 1 + (k-1)\sum_{i} d_i = 1 + 2 \cdot (1+2+4) = 15 \text{ pixels}.
\end{equation}
This $15\times 15$ receptive field enables detection of wide-area degradation (fog, haze, uniform darkness) that single-scale $3\times 3$ models miss.


% ============================================================
% VI. TEMPORAL NC-SSM FOR VIDEO
% ============================================================
\section{Bidirectional Temporal NC-SSM}
\label{sec:temporal}

For video with temporally varying degradation (tunnel entry/exit, flickering lights), frame-independent processing cannot leverage temporal context.
We add a bidirectional temporal NC-SSM after the spatial NC-Conv features:

\subsection{Forward and Backward Paths}
\begin{align}
h_t^{\text{fwd}} &= \sigma_t \cdot \text{Conv1d}_{\text{sel}}^{\text{fwd}}(\mathbf{f})_t \odot \text{gate}(\mathbf{f}_t) \nonumber \\
&\quad + (1-\sigma_t) \cdot \text{Conv1d}_{\text{fixed}}^{\text{fwd}}(\mathbf{f})_t, \\
h_t^{\text{bwd}} &= \text{same structure on reversed } \mathbf{f}, \text{ then re-reversed}.
\end{align}
The two directions are combined via projection: $h_t = \text{FC}([h_t^{\text{fwd}}; h_t^{\text{bwd}}])$.

\subsection{Per-Frame Quality Gate}

A per-frame quality gate $g_t$ controls temporal correction:
\begin{equation}
f_t^{\text{out}} = f_t^{\text{spatial}} + g_t \cdot (h_t^{\text{temporal}} - f_t^{\text{spatial}}),
\end{equation}
with $g_t$ initialized near zero (bias $= -3$, so $\text{sigmoid}(-3) \approx 0.05$), ensuring clean frames are \textbf{protected by default}.

A clean anchor loss penalizes changes to clean frames:
\begin{equation}
\mathcal{L}_{\text{anchor}} = \text{MSE}(f_t^{\text{out}}, f_t^{\text{spatial}}) \quad \text{for clean } t.
\end{equation}

\subsection{Training and Deployment}
\textbf{Training}: Two-phase---pre-train spatial NC-Conv, then freeze spatial and train temporal NC-SSM with anchor loss.
\textbf{Deployment}: Train bidirectional for optimal gradient flow; deploy \textbf{causal} (forward-only) for real-time streaming at 1-frame latency.


% ============================================================
% VII. EXPERIMENTS
% ============================================================
\section{Experiments}
\label{sec:experiments}

\subsection{CIFAR-10-C: Official Corruption Benchmark}
\label{sec:cifar10c}

\paragraph{Setup}
Standard CNN (48--96--192 channels, 251K params) vs.\ NC-Conv (44--88--176, 253K params).
Both trained 80 epochs with AdamW ($\text{lr}=10^{-3}$, weight decay 0.05), warmup 5 epochs, cosine annealing, gradient clipping at 1.0, corruption augmentation (30\%).

\begin{table}[t]
\centering
\caption{CIFAR-10-C benchmark (19 corruptions $\times$ 5 severities, average). Both models use corruption augmentation. Best per-row in \textbf{bold}.}
\label{tab:cifar10c}
\small
\setlength{\tabcolsep}{3pt}
\begin{tabular}{@{}lcc r@{}}
\toprule
\textbf{Corruption} & \textbf{Std CNN} & \textbf{NC-Conv} & $\Delta$ \\
\midrule
Clean (CIFAR-10) & 89.2 & 88.9 & $-0.3$ \\
\midrule
glass\_blur       & 57.3 & \textbf{64.5} & $+7.2$ \\
pixelate          & 65.5 & \textbf{72.4} & $+6.8$ \\
gaussian\_noise   & 75.9 & \textbf{80.1} & $+4.1$ \\
frost             & 76.9 & \textbf{80.0} & $+3.1$ \\
snow              & 74.6 & \textbf{77.4} & $+2.8$ \\
contrast          & 69.1 & \textbf{71.7} & $+2.7$ \\
shot\_noise       & 79.9 & \textbf{82.3} & $+2.4$ \\
speckle\_noise    & 81.0 & \textbf{82.5} & $+1.5$ \\
zoom\_blur        & 73.3 & \textbf{74.3} & $+1.0$ \\
jpeg\_compression & 78.1 & \textbf{78.8} & $+0.7$ \\
motion\_blur      & 71.4 & \textbf{71.9} & $+0.6$ \\
gaussian\_blur    & 71.2 & \textbf{71.7} & $+0.5$ \\
elastic\_transform& 77.2 & \textbf{77.5} & $+0.3$ \\
brightness        & 87.0 & 87.0 & $0.0$ \\
spatter           & \textbf{82.7} & 82.5 & $-0.2$ \\
defocus\_blur     & \textbf{77.5} & 77.3 & $-0.1$ \\
saturate          & \textbf{83.7} & 83.1 & $-0.6$ \\
fog               & \textbf{81.2} & 80.4 & $-0.7$ \\
impulse\_noise    & \textbf{84.1} & 81.7 & $-2.4$ \\
\midrule
\textbf{C10-C Avg} & 76.2 & \textbf{77.7} & $\mathbf{+1.6}$ \\
\bottomrule
\end{tabular}
\end{table}

\begin{table}[t]
\centering
\caption{Severity scaling: NC-Conv's advantage grows monotonically with corruption severity, validating the noise propagation theory (Corollary~\ref{cor:optimal}).}
\label{tab:severity}
\small
\begin{tabular}{@{}cccc@{}}
\toprule
\textbf{Severity} & \textbf{Std CNN} & \textbf{NC-Conv} & $\Delta$ \\
\midrule
1 (mild)   & 84.8 & 85.3 & $+0.5$ \\
2          & 81.5 & 82.4 & $+0.9$ \\
3          & 78.2 & 79.6 & $+1.4$ \\
4          & 72.8 & 74.9 & $+2.1$ \\
5 (severe) & 63.6 & 66.6 & $\mathbf{+2.9}$ \\
\bottomrule
\end{tabular}
\end{table}

NC-Conv outperforms on \textbf{14/19} corruptions ($+1.6\%$ average, Table~\ref{tab:cifar10c}).
The severity scaling (Table~\ref{tab:severity}) shows monotonic increase from $+0.5\%$ at severity 1 to $+2.9\%$ at severity 5---the central prediction of the noise theory.

\subsection{NC-TCN Vision: Dilated Multi-Scale Results}
\label{sec:nctcn_exp}

\begin{table}[t]
\centering
\caption{NC-TCN Vision: dilated multi-scale improves all corruptions over both baselines. All models use corruption augmentation.}
\label{tab:nctcn}
\small
\begin{tabular}{@{}lccc@{}}
\toprule
\textbf{Corruption} & \textbf{Std CNN} & \textbf{NC-Conv} & \textbf{NC-TCN} \\
\midrule
Clean              & 89.0 & 88.8 & 88.8 \\
gaussian\_noise    & 86.3 & 86.5 & \textbf{87.7} \\
brightness\_down   & 86.7 & 86.8 & \textbf{87.5} \\
contrast           & 72.8 & 85.7 & \textbf{88.1} \\
fog                & 70.0 & 85.6 & \textbf{88.2} \\
impulse\_noise     & 76.2 & 76.1 & \textbf{80.8} \\
\midrule
\textbf{Avg Corrupted} & 78.4 & 84.1 & $\mathbf{86.5}$ \\
\bottomrule
\end{tabular}
\end{table}

NC-TCN Vision achieves $+8.1\%$ over Standard CNN ($+2.4\%$ over NC-Conv) on average (Table~\ref{tab:nctcn}).
The dilated receptive field ($15 \times 15$) captures wide-area degradation: fog $+18.2\%$, impulse $+4.6\%$ over Standard CNN.


\subsection{Temporal NC-SSM: Tunnel Video Simulation}
\label{sec:temporal_exp}

\begin{table}[t]
\centering
\caption{Per-frame accuracy on tunnel simulation (8 frames). Quality gate $g_t$: clean$\to$0.004, darkest$\to$0.101---learned without quality labels.}
\label{tab:temporal}
\small
\begin{tabular}{@{}llcccl@{}}
\toprule
& \textbf{Frame} & \textbf{NC-Conv} & \textbf{Bi-SSM} & $\Delta$ & $g_t$ \\
\midrule
$f_0$ & clean    & 89.0 & 87.9 & $-1.1$ & 0.004 \\
$f_2$ & entry    & 68.9 & \textbf{86.7} & $+17.8$ & 0.067 \\
$f_3$ & inside   & 52.6 & \textbf{87.8} & $\mathbf{+35.2}$ & 0.101 \\
$f_4$ & inside+n & 58.6 & \textbf{88.1} & $+29.6$ & 0.092 \\
$f_5$ & exit     & 72.9 & \textbf{84.7} & $+11.8$ & 0.040 \\
$f_7$ & clean    & 89.0 & 88.1 & $-0.9$ & 0.004 \\
\midrule
\multicolumn{2}{l}{Clean ($f_0,f_7$)} & 89.0 & 88.0 & $-1.0$ & -- \\
\multicolumn{2}{l}{\textbf{Degraded} ($f_2$--$f_5$)} & 63.2 & \textbf{86.8} & $\mathbf{+23.6}$ & -- \\
\bottomrule
\end{tabular}
\end{table}

The bidirectional NC-SSM recovers $+35.2\%$ on the darkest frame ($f_3$), with clean frames preserved within $-1.0\%$ (Table~\ref{tab:temporal}).


\subsection{CULane Lane Detection}
\label{sec:culane}

\begin{table}[t]
\centering
\caption{CULane lane detection accuracy (\%) by condition (driver\_37 subset, 3,357 images).}
\label{tab:culane}
\small
\begin{tabular}{@{}lccc@{}}
\toprule
\textbf{Condition} & \textbf{Std CNN} & \textbf{NC-Conv} & $\Delta$ \\
\midrule
Normal & 78.3 & \textbf{84.7} & $+6.5$ \\
Dark   & \textbf{60.0} & 59.0 & $-1.0$ \\
Noise  & 71.6 & \textbf{77.8} & $+6.3$ \\
Fog    & 44.2 & \textbf{65.8} & $\mathbf{+21.7}$ \\
\bottomrule
\end{tabular}
\end{table}

On real driving images (Table~\ref{tab:culane}), NC-Conv achieves $+6.5\%$ on normal and $+21.7\%$ on fog---the most critical adverse condition for tunnel ADAS.


\subsection{Scale Ablation}
\label{sec:scale_exp}

\begin{table}[t]
\centering
\caption{NC-Conv benefit across model scales.}
\label{tab:scale}
\small
\begin{tabular}{@{}lrrcc@{}}
\toprule
\textbf{Scale} & \textbf{Params} & \textbf{Clean} & \textbf{C10-C} & \textbf{Sev.~5} \\
\midrule
1$\times$ Std & 251K & 89.2 & 76.2 & 63.6 \\
1$\times$ NC  & 253K & 88.9 & \textbf{77.7} & \textbf{66.6} \\
\midrule
10$\times$ Std & 1.75M & 91.4 & 79.9 & 68.3 \\
10$\times$ NC  & 1.82M & 91.5 & \textbf{81.2} & \textbf{70.2} \\
\bottomrule
\end{tabular}
\end{table}

NC-Conv provides consistent improvement at both 1$\times$ and 10$\times$ scales (Table~\ref{tab:scale}), with larger benefits at smaller scales where edge deployment requires maximum parameter efficiency.


\subsection{Ablation Studies}
\label{sec:ablation}

\paragraph{$\sigma$ supervision}
Explicit binary supervision hurts by $-4.4\%$.

\paragraph{Per-spatial $\sigma$}
Improves impulse noise from $-2.4\%$ to $-1.3\%$ by enabling local quality detection.

\paragraph{Temporal training}
End-to-end training of spatial+temporal fails; two-phase (pre-train spatial $\to$ freeze $\to$ train temporal) with clean anchor loss resolves this.


% ============================================================
% VIII. NC-VLU: VISION-LANGUAGE EXTENSION
% ============================================================
\section{NC-VLU: Ultra-Lightweight Vision-Language Model}
\label{sec:ncvlu}

Extending the NC framework to vision-language understanding, we propose \textbf{NC-VLU}: the first MCU-scale VLM for \emph{visual keyword spotting}---the vision analog of audio keyword spotting~\cite{choi2026ncssm}.

\subsection{Architecture}
\begin{itemize}
\setlength\itemsep{0pt}
\item \textbf{Vision encoder}: NC-Conv (253K params) $\to$ 128-dim embedding
\item \textbf{Text encoder}: 2-layer Transformer ($d=96$, 4 heads, $\sim$500K params) $\to$ 128-dim embedding
\item \textbf{Alignment}: contrastive InfoNCE loss in shared 128-dim space
\item \textbf{Total}: $\sim$753K params ($<$1\,MB INT8)
\end{itemize}

\subsection{Visual Keyword Spotting for ADAS}
Instead of general-purpose VLM (captioning, VQA), NC-VLU targets a focused set of safety-critical keywords:
``stop sign'', ``pedestrian'', ``traffic light'', ``crosswalk'', etc.
Zero-shot: new keywords can be added without retraining.

\textit{Experimental results to be added.}


% ============================================================
% IX. ANALYSIS
% ============================================================
\section{Analysis}
\label{sec:analysis}

\subsection{Audio-Vision Correspondence}

\begin{table}[t]
\centering
\caption{Audio $\leftrightarrow$ Vision noise-conditioning correspondence.}
\label{tab:mapping}
\small
\begin{tabular}{@{}ll@{}}
\toprule
\textbf{Audio NC-SSM/TCN} & \textbf{Vision NC-Conv/TCN (ours)} \\
\midrule
Selective SSM ($\Delta(x), B(x)$) & Dynamic Conv $\times \alpha(x)$ \\
LTI SSM (fixed) & Static Conv (fixed weights) \\
Per-sub-band SNR $\to \sigma$ & Per-sample/spatial quality $\to \sigma$ \\
$O(\sigma_n^6) \to O(\sigma_n^2)$ & $O(\sigma_n^{4+}) \to O(\sigma_n^2)$ \\
NC-TCN (dilated Conv1d) & NC-TCN Vision (dilated Conv2d) \\
Temporal audio frames & Temporal video frames (Bi-NC-SSM) \\
Audio keyword spotting & Visual keyword spotting (NC-VLU) \\
\bottomrule
\end{tabular}
\end{table}

\subsection{Learned $\sigma$ Gate Behavior}

The spatial $\sigma$ automatically adapts: under normal conditions $\sigma \approx 0.85$ (dynamic path); under dark tunnel conditions $\sigma \approx 0.1$ (static path).
The temporal quality gate tracks degradation exactly: $g = 0.004$ for clean frames, $g = 0.101$ for the darkest frame---all learned without explicit quality supervision.

\subsection{Efficiency and MCU Deployment}
\label{sec:deploy}

\begin{table}[t]
\centering
\caption{Efficiency comparison for edge deployment.}
\label{tab:efficiency}
\small
\begin{tabular}{@{}lrrl@{}}
\toprule
\textbf{Model} & \textbf{Params} & \textbf{INT8} & \textbf{Hardware} \\
\midrule
MobileNetV3-S & 2.5M & 2.5\,MB & CPU (\$30+) \\
UFLD-v2 & 2.7M & 2.7\,MB & CPU (\$30+) \\
\midrule
Std CNN & 251K & 251\,KB & MCU (\$8) \\
\textbf{NC-Conv} & \textbf{253K} & \textbf{253\,KB} & \textbf{MCU (\$8)} \\
+ Temporal SSM & 273K & 273\,KB & MCU (\$8) \\
NC-VLU & 753K & 753\,KB & MCU (\$8) \\
\bottomrule
\end{tabular}
\end{table}

NC-Conv deploys on a \$8 STM32H743 MCU (Cortex-M7, 480\,MHz, 1\,MB SRAM) at 28\,FPS with 35\,ms latency and 0.5\,W power, meeting real-time ADAS requirements.
The temporal NC-SSM adds only 20K parameters; at inference, causal mode enables streaming at 1-frame latency.


% ============================================================
% X. DISCUSSION
% ============================================================
\section{Discussion}
\label{sec:discussion}

\paragraph{Why not just augment more?}
Corruption augmentation (AugMix, DeepAugment) trains the model to handle diverse inputs but applies identical processing to all.
NC-Conv additionally \emph{adapts} its computation based on input quality, providing complementary benefit on top of augmentation.

\paragraph{Impulse noise limitation}
Per-sample $\sigma$ shows $-2.4\%$ on impulse noise where sparse extreme pixels don't alter global statistics.
Per-spatial $\sigma$ partially addresses this ($+1.0\%$), but fully resolving sparse corruption requires pixel-level quality detection---a direction for future work.

\paragraph{Broader applicability}
The NC framework is modality-agnostic: the same $\sigma$-gated dual-path principle applies to audio SSMs~\cite{choi2026ncssm}, audio TCNs~\cite{choi2026nctcn}, spatial CNNs (this work), temporal SSMs (this work), and vision-language models (NC-VLU).
We conjecture it extends to any modality where input-dependent parameterization creates multiplicative noise coupling.


% ============================================================
% XI. CONCLUSION
% ============================================================
\section{Conclusion}
\label{sec:conclusion}

We presented a comprehensive noise-conditioned vision framework transferring audio NC-SSM/TCN theory to visual degradation robustness.
NC-Conv validates the noise propagation theory with monotonic severity scaling on CIFAR-10-C.
NC-TCN Vision achieves $+8.1\%$ through dilated multi-scale quality detection.
The bidirectional temporal NC-SSM recovers $+35.2\%$ on dark tunnel frames while preserving clean accuracy.
On CULane, fog improves by $+21.7\%$.
Quality gates learn degradation profiles without explicit supervision.
At 253K parameters on a \$8 MCU at 28\,FPS, noise-conditioned computation is established as a modality-agnostic framework for degradation-robust edge perception.


% ============================================================
% REFERENCES
% ============================================================
\bibliographystyle{IEEEtran}
\bibliography{refs_tip}

\end{document}
